\section{What ``I have to'' Can Mean}

``I have to.'' In ordinary speech those three words can collapse several claims: what is simply true of me, what act is morally owed, what competent authority has commanded, and how responsible I am for failure. Add force---someone will not let me choose---and the same phrase can make every kind of necessity sound like one pressure. It is not.

\begin{studybox}{Four terms}
\begin{description}[leftmargin=1.45em,style=nextline,itemsep=0.25em,topsep=0.2em]
\item[Dependence] What is true whether or not I consent. A creature receives being; refusal cannot make it uncreated.
\item[Duty] What ought to be chosen. Moral obligation addresses a rational and voluntary power; it binds without physically compelling it.
\item[Precept] A duty, or a determinate mode of fulfilling one, commanded by competent authority for identified subjects and circumstances.
\item[Culpability] Personal responsibility for a concrete failure, assessed by distinguishing matter, knowledge, freedom, intention, and relevant circumstances.
\end{description}
\end{studybox}

The four terms can meet in a single act without collapsing into one another. Their order matters: a precept may specify a duty grounded in a reality it did not create, while culpability still depends upon the agent's concrete knowledge and freedom.

These distinctions prevent several false inferences. ``I depend upon God'' does not by itself yield the canon specifying how the Sunday precept is fulfilled. ``The Church commands this act'' does not mean every external failure is equally culpable. ``This act is owed'' does not mean it is coerced, meritorious without consent, or incapable of love.

\subsection{The layers of law}

Even a genuine duty leaves another question: who may give it common form? Law is an ordinance of reason for the common good, made and promulgated by one who has care of the community (ST I--II, q.~90). Aquinas's categories of eternal, natural, divine, and human law do not name four independent legislatures occupying one plane. Eternal law is divine wisdom ordering acts and ends. Natural law is the rational creature's participation in that order through practical reason. Divine positive law is used here for precepts revealed and promulgated by God within salvation history. Human positive law, civil or ecclesiastical, makes common determinations through competent authority.

The labels keep three sources of normativity distinct: what practical reason knows, what God has revealed or commanded, and what human authority has enacted for common observance. A single practice can stand under more than one description without the sources becoming interchangeable. A theological reason does not by itself enact a canon; a canon does not exhaust the moral, ecclesial, or sacramental good it protects.

\subsection{From general due to common form}

In worship, Aquinas draws the passage precisely. Natural reason dictates that the human person do something in reverence for God. That \emph{this} rather than \emph{that} determinate act be required can be established by divine or human law (ST II--II, q.~81, a.~2, ad 3). Sacrifice likewise belongs generically to natural law, while its concrete signs receive positive determination (q.~85, a.~1, ad 1).

The Lord's Day follows this pattern without becoming a natural-law theorem. Regular outward worship is morally due; its Paschal and juridical determinations come through revelation, apostolic tradition, and ecclesiastical law. Governing law and moral teaching distinguish impossibility and excuse from dispensation and commutation by competent authority. These are juridical distinctions, not private preference.

A concrete obligation must therefore be identified through its authority, promulgated text, jurisdiction, subjects, circumstances, and permitted relief. These questions do not weaken obligation. They identify it. Law states an obligation within a reality it did not originate. A deadline can identify a minimum. It cannot yet show why the day must be shared.
